% Generated draft: requires mathematical and source review.
\documentclass[11pt]{article}
\usepackage[margin=25mm]{geometry}
\usepackage{fontspec}
\setmainfont{Latin Modern Roman}
\usepackage{amsmath,amssymb,mathtools,unicode-math}
\setmathfont{Latin Modern Math}
\usepackage{enumitem,needspace,xurl,hyperref}
\setlength{\parindent}{0pt}
\setlength{\parskip}{5pt}
\setlength{\emergencystretch}{3em}
\newcommand{\N}{\mathbb{N}}
\newcommand{\Z}{\mathbb{Z}}
\newcommand{\Q}{\mathbb{Q}}
\newcommand{\R}{\mathbb{R}}
\newcommand{\C}{\mathbb{C}}
\newcommand{\F}{\mathbb{F}}
\newcommand{\A}{\mathbb{A}}
\newcommand{\PP}{\mathbb{P}}
\newcommand{\E}{\mathbb{E}}
\newcommand{\eps}{\varepsilon}
\newcommand{\dd}{\,\mathrm{d}}
\begin{document}
\section*{UnsolvedMath ID: 7600005}
\textbf{Source problem code:} catalogue entry
\subsection{Definitions and mathematical statement}
For a first-order theory $T$ and an infinite cardinal $\lambda$, let
\[
I(T,\lambda)=|\{M\models T:|M|=\lambda\}/\cong|
\]
be the model spectrum. An $\aleph_1$-saturated model realizes every consistent type over each countable parameter set. Let $I_{\aleph_1\text{-}\mathrm{sat}}(T,\lambda)$ count the isomorphism types of such models of size $\lambda$.

An abstract elementary class (AEC) $(K,\preccurlyeq)$ has a fixed vocabulary and a coherent strong-substructure relation, is closed under isomorphisms and directed unions, and satisfies a L\"owenheim--Skolem bound. Its spectrum $I(K,\lambda)$ similarly counts isomorphism types of members of size $\lambda$.

Formulate and prove the intended main-gap dichotomies in each of the source's settings:
\begin{enumerate}
\item Uncountable first-order theories.
\item Abstract elementary classes.
\item $\aleph_1$-saturated models of a countable theory.
\end{enumerate}
A main-gap statement seeks a sharp division between a structural/small-spectrum side and a nonstructure/large-spectrum side, with the correct cardinal ranges and quantitative bounds. The source gives none of those bounds or hypotheses. In particular, uncountable theory size must not be confused with merely studying uncountable models of a countable theory. For AECs, no amalgamation, tameness, cardinal-arithmetic, or other additional condition is supplied.

The accompanying summary reports Shelah's original main-gap theorem for the countable first-order setting as proved, while retaining the listed extensions as broader unresolved programs. That result is background and is not used to claim a universal arbitrary-AEC dichotomy.
\subsection{Short English statement}
Extend main-gap classification beyond the solved countable first-order setting, including uncountable theories, AECs, and $\aleph_1$-saturated models of countable theories. The exact spectrum bounds, hypotheses, and cardinal ranges are missing.
\subsection{Sources}
\begin{itemize}
\item Wikipedia contributors, \emph{List of unsolved problems in mathematics}, ``Model theory and formal languages,'' bullet 5, pinned revision 1366636767, \url{https://en.wikipedia.org/w/index.php?title=List_of_unsolved_problems_in_mathematics&oldid=1366636767#Model_theory_and_formal_languages}; accessed 2026-07-29, MediaWiki source extraction. Wikipedia text is CC BY-SA 4.0, with attribution/share-alike handling required for release.
\item Shelah's countable first-order main-gap theorem is named in the supplied research summary, without a primary publication citation or exact statement of the intended extensions.
\item Local registry \texttt{data/unsolvedmath/problems.json}, permanent UnsolvedMath ID $7600005$, source code catalogue entry, \emph{The main gap conjecture}. Dataset attribution: Ulam AI and the UnsolvedMath Contributors.
\end{itemize}
\end{document}
